\chapter[Conservation]{Covariant Conservation of the Total Stress Tensor}
\label{ch:conservation}

\paragraph{Learning outcomes.}
After this chapter the reader should be able to:
\begin{enumerate}
  \item define \(T_{\mu\nu}^{\mathrm{total}}\) as the sum of matter,
        tensor-block, and entropic contributions;
  \item obtain \(\Tent_{\mu\nu}\) by metric variation of
        \(U_{\mathrm{ent}}\);
  \item prove that diffeomorphism invariance forces
        \(\nabla^\mu T_{\mu\nu}^{\mathrm{total}}=0\) on shell;
  \item reduce that identity to a Newtonian force balance and see
        which drags are allowed;
  \item enter Chapter~12 knowing that vacuum energy is a
        conservation problem as well as a cosmological-constant
        problem.
\end{enumerate}

Chapter~10 made \(\nabla_\mu\) the Levi-Civita connection of the
composite metric. This chapter asks what that connection is allowed
to act on.

\section{The split}
\label{sec:ch11-split}

The infrared field equation contemplated from here on is
\begin{equation}
  \label{eq:EFE-total}
  G_{\mu\nu}+\Lambda_{\mathrm{eff}}g_{\mu\nu}
  =8\pi\Geff\,T_{\mu\nu}^{\mathrm{total}},
\end{equation}
with
\begin{equation}
  \label{eq:T-total-def}
  T_{\mu\nu}^{\mathrm{total}}
  =T_{\mu\nu}^{\mathrm{matter}}
  +T_{\mu\nu}^{\mathrm{tensor}}
  +\Tent_{\mu\nu}.
\end{equation}

\paragraph{Matter.}
\(T_{\mu\nu}^{\mathrm{matter}}\) is the metric variation
\eqref{eq:T-def} of every action for fields that are not the
condensate itself: ordinary fluids, electromagnetism, baryons. On
shell it satisfies \(\nabla^\mu T_{\mu\nu}^{\mathrm{matter}}=0\)
\emph{only} if those fields do not exchange energy with the
condensate. In NAMCG they may. The separate conservation of matter
is therefore not assumed.

\paragraph{Tensor block.}
\(T_{\mu\nu}^{\mathrm{tensor}}\) is the metric stress of the
quadratic action~\eqref{eq:EH-lin}. After Theorem~\ref{thm:mT} it
is the linearized gravitational-wave stress, of the Isaacson type
when averaged over wavelengths short compared with the background
curvature. It is kept in~\eqref{eq:T-total-def} for bookkeeping and
is negligible in the static Newtonian laboratories of
Chapters~14--15.

\paragraph{Entropic piece.}
\(\Tent_{\mu\nu}\) is defined by the same variational rule applied
to the mutual-information potential of Definition~\ref{def:Uent},
promoted to a functional of the metric:
\begin{equation}
  \label{eq:Tent-var}
  \Tent_{\mu\nu}
  =-\frac{2}{\sqrt{-g}}
    \frac{\delta S_{\mathrm{ent}}[g]}{\delta g^{\mu\nu}},
  \qquad
  S_{\mathrm{ent}}[g]
  =-\int\mathrm{d}t\,U_{\mathrm{ent}}[g].
\end{equation}
The promotion \(U_{\mathrm{ent}}\to S_{\mathrm{ent}}[g]\) uses the
fact that every entropy \(S(A)\) in~\eqref{eq:Uent-def} is computed
from a reduced state on a region whose causal structure depends on
\(g_{\mu\nu}\). Varying the metric varies the regions and varies
the modular Hamiltonians.

\section{Noether's second theorem, applied once}
\label{sec:ch11-noether}

Let \(S_{\mathrm{IR}}[g,\Psi]\) be the infrared action whose
Euler--Lagrange equation for \(g_{\mu\nu}\) is~\eqref{eq:EFE-total}:
Einstein--Hilbert plus \(\Lambda_{\mathrm{eff}}\), plus matter, plus
\(S_{\mathrm{ent}}\), plus the tensor-block action. Diffeomorphism
invariance states that for every compactly supported vector field
\(\xi^\mu\),
\begin{equation}
  \label{eq:diff-inv}
  \delta_\xi S_{\mathrm{IR}}=0
\end{equation}
identically, off shell as a statement about the integrand up to a
boundary term. Noether's second theorem then supplies
\begin{equation}
  \label{eq:noether-2}
  \nabla_\mu
  \Bigl(
    \frac{1}{8\pi\Geff}
    \bigl(G^{\mu\nu}+\Lambda_{\mathrm{eff}}g^{\mu\nu}\bigr)
    -T^{\mu\nu}_{\mathrm{total}}
  \Bigr)
  =0
\end{equation}
as an identity, once the matter and condensate fields \(\Psi\) are
on shell. The contracted Bianchi identity
\(\nabla_\mu G^{\mu\nu}=0\) and metric compatibility
\(\nabla_\mu g^{\mu\nu}=0\) reduce~\eqref{eq:noether-2} to the
on-shell conservation law
\begin{equation}
  \label{eq:cons-law}
  \nabla_\mu T^{\mu\nu}_{\mathrm{total}}=0.
\end{equation}

\begin{theorem}[Conservation of the total stress]
  \label{thm:conservation}
  Assume:
  \begin{enumerate}
    \item the infrared action \(S_{\mathrm{IR}}\) is a scalar under
          diffeomorphisms of the composite metric, including the
          entropic functional \(S_{\mathrm{ent}}[g]\);
    \item the Einstein equation is~\eqref{eq:EFE-total};
    \item matter fields and condensate fluctuations are on shell
          with respect to \(S_{\mathrm{IR}}\).
  \end{enumerate}
  Then
  \begin{equation}
    \label{eq:cons-thm}
    \nabla^\mu
    \bigl(
      T_{\mu\nu}^{\mathrm{matter}}
      +T_{\mu\nu}^{\mathrm{tensor}}
      +\Tent_{\mu\nu}
    \bigr)
    =0.
  \end{equation}
  Separate conservation of any one summand is not implied.
\end{theorem}

The Einstein-plus-extra-tensor form of the infrared equation
is collected in Appendix~\ref{app:field-eq}.

\begin{proof}
  Diffeomorphism invariance of \(S_{\mathrm{IR}}\) is
  Assumption~1. Noether's second theorem yields~\eqref{eq:noether-2}.
  Bianchi and~\eqref{eq:EFE-total} yield~\eqref{eq:cons-thm}. The
  last sentence is the observation that Noether's theorem constrains
  only the sum that couples to \(g_{\mu\nu}\).
\end{proof}

\begin{remark}[The real hypothesis]
  Everything in the proof is standard except Assumption~1 applied
  to \(S_{\mathrm{ent}}\). Mutual information~\eqref{eq:MI} is a
  scalar for each fixed partition of a spatial slice. Promoting it
  to a diffeomorphism scalar on a dynamical spacetime requires that
  the assignment
  \((g,A,B)\mapsto I(A:B)\) commute with Lie dragging of the
  regions. That is true if \(A\) and \(B\) are defined
  geometrically (causal domains, outermost trapped surfaces,
  constant-mean-curvature balls of specified volume) and fails if
  they are defined by coordinate inequalities such as
  \(r\le R\) in a preferred chart. The book uses geometric
  regions.
\end{remark}

\section{Structure of the entropic stress}
\label{sec:ch11-Tent}
\label{sec:ch11-Tent}

A first variation of \(S_{\mathrm{ent}}\) on a static background
produces a conservative extra potential. A first variation on a
time-dependent background produces, in addition, a piece odd under
velocities. The split
\begin{equation}
  \label{eq:Tent-split}
  T_{\mu\nu}^{\mathrm{ent}}
  =T_{\mu\nu}^{\mathrm{ent,rev}}
   +T_{\mu\nu}^{\mathrm{ent,irrev}}
\end{equation}
is the covariant version of ``energy from mutual information'' plus
``drag from changing mutual information.''

On two worldlines \(\gamma_A,\gamma_B\) in a nearly Minkowski
region one has, at leading post-Newtonian order,
\begin{equation}
  \label{eq:U-PN}
  U_{\mathrm{ent}}
  =-\lambda I\bigl(\ell_A,\ell_B,d(\gamma_A,\gamma_B)\bigr),
\end{equation}
and the four-force on body \(A\) is
\begin{equation}
  \label{eq:four-force}
  f^\mu_A
  =-\frac{D}{D\tau_A}
    \Bigl(
      \frac{\partial L_{\mathrm{ent}}}{\partial \dot x_{A\mu}}
    \Bigr)
    +\frac{\partial L_{\mathrm{ent}}}{\partial x_{A\mu}},
\end{equation}
with \(L_{\mathrm{ent}}=-U_{\mathrm{ent}}\) along the worldlines.
The term built from \(\partial I/\partial d\) at vanishing velocity
is reversible and conservative. The term built from
\(\partial I/\partial\dot d\) or from the time derivative of
\(\ell_A,\ell_B\) is irreversible in the thermodynamic sense of
Chapter~6: it tracks a change of entanglement across the motion,
not a stored potential.

The Newtonian reduction of the irreversible piece, for a reduced
mass \(\mu\) and a separation \(d\), is of the form
\begin{equation}
  \label{eq:airrev}
  \mathbf{a}_{\mathrm{irrev}}
  =-\eta(d,\ell_A,\ell_B)\,
    \frac{\mathbf{v}}{d^{n}},
\end{equation}
with a positive coefficient \(\eta\) if mutual information decreases
under adiabatic separation of two thermal cores. The source notes
use \(n=4\) as a working exponent compatible with a second
derivative of a \(1/d^2\) mutual-information tail
\eqref{eq:I-falloff}. Chapter~15 adopts
\begin{equation}
  \label{eq:airrev-lab}
  \mathbf{F}_{\mathrm{diss}}
  =-\eta\,\frac{\mu}{d^4}\,\mathbf{v}
\end{equation}
as a laboratory, not as a derived multipole formula.

\begin{proposition}[Allowed Newtonian drags]
  \label{prop:allowed-drag}
  A two-body equation
  \begin{equation}
    \label{eq:newton-drag}
    \mu\ddot{\mathbf{r}}
    =-\nabla U_{\mathrm{N}}-\nabla U_{\mathrm{ent}}^{\mathrm{(rev)}}
     +\mathbf{F}_{\mathrm{diss}}
  \end{equation}
  is the reduction of a covariantly conserved \(\Tent_{\mu\nu}\)
  only if there exists an energy-balance identity
  \begin{equation}
    \label{eq:energy-balance}
    \frac{\mathrm{d}}{\mathrm{d}t}
    \bigl(E_{\mathrm{mech}}+E_{\mathrm{ent}}\bigr)
    =\mathbf{F}_{\mathrm{diss}}\cdot\mathbf{v}
     +O(c^{-2})
  \end{equation}
  with \(E_{\mathrm{ent}}\) built from \(I(A:B)\), and only if the
  equal-and-opposite force on the second body holds up to radiation
  reaction of Isaacson type. A drag that dissipates mechanical
  energy into nowhere, with no \(E_{\mathrm{ent}}\) and no
  gravitational radiation, is not allowed.
\end{proposition}

The laboratory~\eqref{eq:airrev-lab} satisfies
Proposition~\ref{prop:allowed-drag} if and only if the integrator
tracks an entanglement energy whose decrease equals
\(-\mathbf{F}_{\mathrm{diss}}\cdot\mathbf{v}\). Chapter~15 will do
that bookkeeping explicitly.

\section{Exchange with matter}
\label{sec:ch11-exchange}

Equation~\eqref{eq:cons-thm} permits
\begin{equation}
  \label{eq:exchange}
  \nabla^\mu T_{\mu\nu}^{\mathrm{matter}}
  =-\nabla^\mu\Tent_{\mu\nu}
   -\nabla^\mu T_{\mu\nu}^{\mathrm{tensor}}.
\end{equation}
Matter may heat or cool because entanglement of the condensate
cores is changing. In a tidal disruption event the geometric
reading is: as the disrupted star is sheared, the mutual information
between debris parcels and the compact core changes, and the
work done by \(\mathbf{F}_{\mathrm{diss}}\) is accounted in
\(E_{\mathrm{ent}}\), not lost.

This is still not a model of microphysical dissipation in a fluid
with a viscosity measured in the laboratory. It is a channel
required by conservation once \(U_{\mathrm{ent}}\) is admitted as
part of the action.

\section{Vacuum energy as a conservation constraint}
\label{sec:ch11-vacuum}

A term \(\rho_\Lambda g_{\mu\nu}\) with constant \(\rho_\Lambda\) is
conserved by itself:
\begin{equation}
  \label{eq:Lambda-cons}
  \nabla^\mu(\rho_\Lambda g_{\mu\nu})=0.
\end{equation}
A term whose prefactor varies, \(\rho_\Lambda(x)g_{\mu\nu}\), is
not conserved unless
\begin{equation}
  \label{eq:Lambda-var}
  \partial_\mu\rho_\Lambda=0
\end{equation}
or additional current terms are supplied. The heat-kernel
density~\eqref{eq:W-Lambda} is a constant in the approximation that
\(\LUV\) and \(\mathrm{Tr}(\mathbb{I})\) are constant. If Chapter~12
regulates \(\rho_\Lambda\) by a boundary functional that depends on
the metric induced on \(\partial M\), the regulated object
\(\Lambda_{\mathrm{eff}}[g]\) must still obey~\eqref{eq:cons-thm}.
That requirement is a design constraint on HBEV, not an optional
aesthetic.

\section{What this chapter has not done}
\label{sec:ch11-not}

Theorem~\ref{thm:conservation} does not compute \(\eta\), does not
fix \(n=4\), does not evaluate \(I(A:B)\) for Kerr exteriors, and
does not replace quadrupole radiation. Gravitational waves remain
the dissipative channel of the tensor block. Entropic drag is an
additional channel, allowed only because it sits inside a conserved
total tensor.

The theorem also does not justify writing an arbitrary function
\(\mathbf{a}_{\mathrm{irrev}}(\mathbf{r},\mathbf{v})\) and calling it
NAMCG. Phenomenological chapters that introduce a drag must exhibit
the \(E_{\mathrm{ent}}\) of Proposition~\ref{prop:allowed-drag} or
must state that they are studying a reduced laboratory.

\section{Exercises}
\label{sec:ch11-exercises}

\begin{enumerate}
  \item \textbf{Bianchi reminder.}
    From \(\nabla^\mu G_{\mu\nu}=0\) and~\eqref{eq:EFE-total},
    obtain~\eqref{eq:cons-thm} in three lines.
  \item \textbf{Two-body check.}
    For a potential \(U(d)\) depending only on separation, show
    that the Newtonian forces are equal and opposite. Which
    symmetry of \(I(A:B)\) is the relativistic ancestor of that
    statement?
  \item \textbf{Energy bookkeeping.}
    Take~\eqref{eq:airrev-lab} with constant \(\eta\) and compute
    \(\mathbf{F}_{\mathrm{diss}}\cdot\mathbf{v}\). Propose an
    \(E_{\mathrm{ent}}(d,\dot d)\) whose time derivative matches
    that power. Is your \(E_{\mathrm{ent}}\) unique?
  \item \textbf{Illegal drag.}
    Consider \(\mathbf{F}=-\gamma\mathbf{v}\) with constant
    \(\gamma\), independent of \(d\). Why does this fail
    Proposition~\ref{prop:allowed-drag} for two bodies that can
    be separated to infinity at finite cost in \(I(A:B)\)?
  \item \textbf{Variable vacuum energy.}
    Let \(\rho_\Lambda=\rho_\Lambda(\phi)\) for a scalar
    \(\phi\). Construct the additional stress that restores
    conservation. Compare with a potential \(V(\phi)\) in
    ordinary scalar-field cosmology.
  \item \textbf{Isaacson average.}
    Write the Isaacson stress of a high-frequency TT wave and
    confirm that its divergence vanishes after a short-wavelength
    average. That is the tensor-block contribution to
    \eqref{eq:cons-thm} in a radiative spacetime.
  \item \textbf{Coordinate regions.}
    Give an explicit example of two coordinate-defined regions
    whose mutual information is \emph{not} a diffeomorphism
    scalar. This is the counterexample that forces the remark
    after Theorem~\ref{thm:conservation}.
\end{enumerate}

\section{Bridge to Chapter~12}
\label{sec:ch11-bridge}

Conservation is now in force for whatever vacuum energy Chapter~8
left untamed. A constant \(\rho_\Lambda\) is conserved and
disastrous. A regulated \(\Lambda_{\mathrm{eff}}\) must remain
conserved after the regulation. Chapter~12 introduces the
Holographic Boundary-Entropy Vacuum: a boundary functional built
from \(K/8\pi\Geff\) and an entropy density
\(S_{\mathrm{ent}}(h_{ij},\phi)\) whose stationarity sets
\(\Lambda_{\mathrm{eff}}\) by an infrared entanglement datum
rather than by \(\LUV^4\). The design constraint just stated---that
the resulting \(\Lambda_{\mathrm{eff}}g_{\mu\nu}\) still sit
inside Theorem~\ref{thm:conservation}---is the first test that
construction must pass.